% from ega4-2-fr.tex
\oldpage[IV-2]{181}
\begin{proposition}[6.15.10]
\label{IV.2.s6.15.10-fr}
Soient $Y$, $Y'$ deux préschémas localement noethériens, $g:Y'\to Y$ un
morphisme normal (6.8.1), $f:X\to Y$ un morphisme; posons
$X'=X\times_Y Y'$ et soient $p:X'\to X$ la projection canonique, $x'$ un
point de $X'$. Alors, si $X$ est réduit (resp. géométriquement unibranche,
resp. intègre et géométriquement unibranche) au point $x=p(x')$, $X'$ est
réduit (resp. géométriquement unibranche, resp. intègre et
géométriquement unibranche) au point $x'$.
\end{proposition}

\begin{proof}
En vertu de (I, 3.6.5), on peut se borner au cas où
$Y=\operatorname{Spec}(A)$, $Y'=\operatorname{Spec}(A')$,
$X=\operatorname{Spec}(B)$, où $A$, $A'$ et $B$ sont des anneaux locaux,
les homomorphismes $A\to A'$ et $A\to B$ étant locaux, $A$, $A'$ des
anneaux noethériens, et $x$ étant le point fermé de $X$; il s'agit de
prouver que si $B$ est réduit (resp. géométriquement unibranche, resp.
intègre et géométriquement unibranche) alors $B'=B\otimes_A A'$ est réduit
(resp. $\operatorname{Spec}(B')$ est géométriquement unibranche aux points
de $p^{-1}(x)$, resp. $\operatorname{Spec}(B')$ est intègre et
géométriquement unibranche en ces points). Supposons d'abord seulement
$B$ réduit; $B$ est réunion de la famille croissante de ses
sous-$A$-algèbres de type fini $B_\lambda$, et par platitude, $B'$ est
réunion de la famille filtrante croissante de ses sous-$A'$-algèbres
$B'_\lambda=B_\lambda\otimes_A A'$. Or, le morphisme
$\operatorname{Spec}(B'_\lambda)\to\operatorname{Spec}(B_\lambda)$ est
normal (6.8.2), et a fortiori réduit; le fait que $B'_\lambda$ soit réduit
est donc conséquence de (3.3.5); on conclut que $B'$ lui-même est réduit par
(5.13.2).

Supposons maintenant $B$ géométriquement unibranche; compte tenu de (I,
5.1.8), on peut en outre supposer $B$ réduit, donc intègre. Soit $C$ sa
clôture intégrale, qui est par hypothèse un anneau local; posons
$Z=\operatorname{Spec}(C)$, $C'=C\otimes_A A'$,
$Z'=\operatorname{Spec}(C')=Z\times_X X'$, de sorte que l'on a le diagramme
commutatif
\[
\xymatrix{
Z \ar[d]_l & Z' \ar[l] \ar[d]_{l'}\\
X & X' \ar[l]_p
}
\]
En vertu de la première partie du raisonnement, $X'$ est réduit; d'autre
part, comme $Z'=Z\times_Y Y'$, il résulte de (6.14.1) que $Z'$ est normal
(et a fortiori géométriquement unibranche). Comme $f$ est entier et
birationnel, il en est de même de $f'$ par (6.15.4.1), puisque $p$ est
plat. Enfin, l'hypothèse que $X$ est géométriquement unibranche au point
$x$ entraîne que $f$ est radiciel en ce point (6.15.5); on conclut donc de
(6.15.3.1) que $f'$ est radiciel en tout point $x'\in p^{-1}(x)$, et il
résulte alors de (6.15.5) que $X'$ est géométriquement unibranche (donc
intègre puisqu'il est réduit) en chacun de ces points.
\end{proof}

\begin{remarks}[6.15.11]
\label{IV.2.s6.15.11-fr}
(i) Dans la démonstration de (6.15.10), on ne peut se dispenser de faire
appel à (6.14.1), même lorsque $X=Y$, car on fait intervenir la clôture
intégrale de l'anneau $B$, qui n'est pas nécessairement un anneau
noethérien même si $B$ l'est.

(ii) L'exemple (6.5.5, (ii)) montre que dans (6.15.10), on ne peut remplacer
« géométriquement unibranche » par « intègre » dans l'énoncé, même si
le corps résiduel $k(x)$ est algébriquement clos et le morphisme $g$ étale.

